\section{Introduction}
\label{sec:intro}

Reported earnings at the largest AI spenders now hinge on one management estimate: how many years
a server lasts. Between 2020 and 2025 Microsoft, Alphabet, Amazon, Meta, Oracle and CoreWeave
lengthened the useful lives of their servers, and Amazon partly reversed course for a subset of
them in 2025 (Section~\ref{sec:background} and Table~\ref{tab:lives}). Whether five- or six-year
straight-line lives are defensible for graphics processing units (GPUs) bought at the peak of an
investment boom is now a live question for investors, auditors and regulators.

The first academic and practitioner estimates of accelerator depreciation have appeared, and they
disagree. \citet{amissahocran2026collateral} fit a rent--age profile at one date and find a
retention path consistent with roughly 40\% a year; \citet{chincholikar2026economic} use executed
resale prices and find 24--27\%; \citet{huang2026gpu} measure rental contribution on a
peer-to-peer marketplace and find 17--32\%; and archived posted on-demand prices show declines
from roughly 15\% a year to none at all \citep{ccir2026rental}. These studies measure different
things: resale value, net cash flow or posted rent, at one date or along a path.

This paper puts those measurements on one basis, using the object that is easiest to observe
over time: what it costs to rent an accelerator for an hour. I build a panel of \nPanelRows{}
listing-months of posted rents for seven NVIDIA chips, from the V100 (2017) to the B200 (2024),
covering October 2017 to November 2026. It combines Amazon Web Services (AWS) price-list history,
archived pricing pages, the SkyPilot catalog history and live collection, across hyperscalers and
neoclouds and across on-demand, spot and reserved contracts (Section~\ref{sec:data}).

The central econometric problem is old. Age equals calendar time minus vintage, so with chip and
month fixed effects the linear part of the age profile is not identified \citep{hall1968technical}.
Following the advice to report each chip's realized path first, I then run four identification
strategies side by side: a quality-adjusted vintage effect, a frontier-relative obsolescence term
after \citet{geske2007computers}, a successor event study, and a shape-only specification that
sweeps an assumed linear rate (Section~\ref{sec:method}). Inference is at the chip level, which
leaves seven clusters.

Three results stand out.

\emph{Realized rents fall slowly.} Holding provider, region and contract fixed, the V100's rent
changes by \nPathVHundredFour{} log points over its first four years and \nPathVHundredNine{}
over nine; the A100 40GB's falls to \nPathAHundredFive{} by year five. The H100's rent rose by
\nPathHHundredOne{} in its first year and was \nPathHHundredThree{} at three years, before 2026's
shortage brought it back to \nPathHHundredFour{} at four years.

\emph{The shape of the age profile depends on the market, and the data mostly measure pricing
habits.} \nListUnchanged{} of month-to-month on-demand and reserved transitions in the panel show
no price change; AWS has charged \$\nAwsVHundredPrice{} an hour for a V100 on demand in every year
from \nAwsVHundredFirstYear{} to \nAwsVHundredLastYear{}. Spot prices give a back-loaded profile
(rate over years 0--2 minus the rate over years 2--4 of \nFourShapeSpot), while list prices that
do change give a front-loaded one (\nShapeListMovers, provider-clustered standard error
\nShapeListMoversSe). \ajcheck{Which of these is ``the'' age profile is a choice of market, not a
technical detail.}

\emph{With valid inference, the data do not pin down the rate or the shape.} Under the
unrestricted wild cluster bootstrap that the analysis pre-registered, the pooled shape estimate
(\nFourShape) excludes zero. A Monte Carlo on this panel's design shows that the unrestricted
bootstrap's 95\% intervals cover the truth only \nCovWcuLo--\nCovWcuHi{} of the time, while the
restricted version covers about \nCovWcrHi; under the restricted bootstrap the shape has
$p = \nFourShapeWcrP$ and a 95\% interval of [\nFourShapeWcrLo, \nFourShapeWcrHi]. No
identification strategy that identifies a linear rate finds a significant decline over years 0--4.

Turning rents into asset values (Section~\ref{sec:values}), the present value of net rent falls
along the realized path by \nValRealLostThree{} by year three for an H100, which is slower than any
five- or six-year straight-line schedule. The H100's neocloud rent of \$\nValRentNow{} an hour
compares with roughly \$\nValOpex{} of power and maintenance, so rents never fall to operating cost within the
nine-year horizon: economic life is set by obsolescence and power scarcity, not by rents reaching
operating cost. A separate benchmark, term-structure and retirement analysis
(Section~\ref{sec:triangulation}) places the power-constrained retirement age at
\nPThreeMedian{} years and the life to operating cost at \nPTwoLopexLo--\nPTwoLopexHi{} years,
with book lives in between.

\ajcheck{The contribution is a reconciliation. The spread of published rates comes mostly from
three choices that are rarely stated: which price series is read (resale, spot, posted list, or
transacted rent), which identification assumption pins the linear age term, and whether value is
taken to end when rent reaches operating cost or when a better chip displaces the old one in a
power-constrained slot.} The paper's open panel, code and identification range let readers make
those choices explicitly.

\paragraph{Related literature.} The paper applies the vintage-price approach to depreciation
\citep{hulten1981estimation, oliner1993constant, doms2004computers, geske2007computers} and
hedonic price measurement \citep{rosen1974hedonic, berndt1993price, triplett2006handbook} to a new
asset class with an active rental market, as \citet{gavazza2011leasing} does for aircraft. Its
identification problem is the age--period--cohort problem of \citet{hall1968technical} and
\citet{deaton1997analysis}. On inference with few clusters it follows
\citet{cameron2008bootstrap}, \citet{mackinnon2017wild}, \citet{webb2023reworking} and
\citet{mackinnon2023cluster}. On the accounting side it relates to \citet{keating1999depreciation}.
Recent work measures cloud prices \citep{byrne2018cloud, sawyer2023ppi} and GPU price-performance
\citep{hobbhahn2022trends}, and a fast-growing practitioner and working-paper literature argues
about accelerator lives \citep{dutta2026priced, rosenzweig2026second}.
